\chapter{Binomial Distribution}
\lecture{2}{28 Apr. 21:51}{Second Lecture}
\label{chap:binomial-distribution}

\section{Binomial Expansion}
\label{sec:binomial-expansion}
\[
  (a+b)^{n} = \sum_{r=0}^{n}{\binom{n}{r} a^{n-r}b^{r}}
.\] 

\subsection{Properties of Binomial Expansion}
\label{ssec:properties-of-binomial-expansion}

\begin{enumerate}
  \item There are \(n+1\) elements.
  \item \(\binom{n}{r} = \binom{n}{n-r}\)
  \item \((n-r)+r = n\) 
  \item \(t_{r+1} = \binom{n}{r} q^{n-r} p^{r}\)
\end{enumerate}

\section{Head Tail Experiment(HTE)}
\label{sec:head-tail-experiment}

\subsection{Outcomes of \hyperref[sec:head-tail-experiment]{HTE}}
\label{ssec:outcomes-of-hte}

A HTE has 2 possible outcomes. So it is
\href{sec:bernoulli-process}{Bernoulli Process}.
\begin{enumerate}
  \item Head (H)
  \item Tail (T)
\end{enumerate}

We set Head to be success (1) and Tail to be failure (0).\\

Then, we have the probability of success as:
\[p = \frac{1}{2}\] 
And of failure as:
\[q = \frac{1}{2}\] 

\textbf{Experiment:}\\
Let, \(n\) be the \(\#\)trials i.e we toss a coin \(n\)
times.\\[10pt]
The outcome of this experiment is that we get \(r\) \(\#\)
success and \((n-r)\) failure.\\

\begin{note}
  The Probability that we get exactly  \(r\) \(\#\) success
  is same as:
  \[
    \binom{n}{r}q^{n-r}p^{r}
  \]
  because as stated before, there are \(r\) success. So the
  probability:\\
  \marginpar{\(p\) is fixed.}
  \[P(H \cap H \cap H \ldots \text{r times}) = \prod^{r}p=p^{r}\]
  and similarly, \(q^{n-r}\) i.e.

  \begin{equation}
    \label{eqn:binomial-distribution}
    \begin{displaystyle}
      P_X(r) = P(X=r) = \binom{n}{r}q^{n-r}p^{r}
    \end{displaystyle}
    \marginpar{ See
      \hyperref[def:probability-mass-function]{Probability
    Mass Function}.}
  \end{equation}

\end{note}

\begin{note}
  The Probability that we get at most  \(r\) \(\#\) success
  is same as \(F_X(r)\):

  \begin{equation}
    \label{eqn:cdf-binomial-distribution}
    \begin{displaystyle}
      F_X(r) = P(X\le r) = \sum_{n=0}^{r} P_X(r) 
    \end{displaystyle}
    \marginpar{ See
      \hyperref[def:cumulative-distribution-function]{Cumulative
    Distribution Function}.}
  \end{equation}
\end{note}

\section{Exclusive and Independent Events}
\label{sec:exclusive-and-independent-events}

\begin{definition}[Mutually Exclusive Events]
  \label{def:mutually-exclusive-events}
  Two events \( A \) and \( B \) are mutually exclusive
  (disjoint) if they cannot occur simultaneously.
  Mathematically:
  \[
    P(A \cap B) = 0
  \]
\end{definition}

\begin{definition}[Independent Events]
  \label{def:independent-events}
  Two events \( A \) and \( B \) are independent if the
  occurrence of one does not affect the probability of the
  other. Mathematically:
  \[
    P(A \cap B) = P(A) \cdot P(B)
  \] or,
  \[
    P(A \mid B) = P(A)
  \]
\end{definition}

\section{Bernoulli Process}
\label{sec:bernoulli-process}
In this whole experiment, we have tossed a coin many times.
First we fixed how many times we toss. (We fix \(n\)). We
have the two possible events, one head and the other tail.
These events are mutually exclusive and independent. We also
know the probability of success and this probability is
constant (fixed) throughout the experiment.

What we have defined here is a process called
\textbf{Bernoulli Process}.

\begin{remark}
  A Bernoulli process consists of many independent \emph{Bernoulli Trials},
  each trial having exactly two possible outcomes (often
  called success and failure).
\end{remark}

\begin{definition}[Bernoulli Process]
  \label{def:bernoulli-process}
  A \emph{Bernoulli process} is a sequence of independent
  and identically distributed (i.i.d.) \emph{Bernoulli trials},
  where each trial results in one of two possible outcomes:
  "success" (with probability \(p\)) or "failure" (with
  probability \(1-p\)). The number of trials \(n\) is fixed
  in advance, and the probability of success \(p\) remains
  constant throughout the experiment.

  Formally, the process consists of random variables \(X_1,
  X_2, \ldots, X_n\), where each
  \[
    R_{X_i} = 
    \begin{cases}
      1, & \text{with probability } p \quad (\text{success}) \\
      0, & \text{with probability } 1-p \quad (\text{failure})
    \end{cases}
  \]
  represents range of random variable for \(i^{th}\) trial among the \(n\)
  trials and \(X_i\) are independent.\\

  Here,
  \begin{enumerate}
    \item \(n\) is fixed
    \item Events are mutually exclusive
    \item Trials are independent
    \item \(p\) is known and fixed throughout the
      experiment.
  \end{enumerate}
\end{definition}


\section{Probability Distribution for Bernoulli Trial}
\label{sec:bernoulli-distribution}
The distribution resulted by \hyperref[eqn:binomial-distribution]{this probability mass
function} for \(n=1\).
\[ P_X(r) = p^{r} \] Here, \(r\in \{0,1\} \) since there cannot be more than 1
success in one trial. It is a special case of \href{Binomial
Distribution}{sec:binomial-distribution}

\begin{enumerate}
  \item \(\mu=p\)
  \item \(\sigma^{2}=p(1-p)=pq\)
\end{enumerate}

\section{Binomial Distribution}
\label{sec:binomial-distribution}
The distribution resulted by
\hyperref[eqn:binomial-distribution]{this probability mass
function} is a binomial distribution.

\subsection{Symmetricity of Binomial Distribution}
\label{ssec:symmetricity-of-binomial-distribution}

\begin{remark}
  If \(p=0.5\), distribution is symmetric\\
  \begin{explanation}
    If probability of success is same as the failure, there
    is no skewness in the distribution.\\
  \end{explanation}

  If \(p<0.5\), distribution is skewed right(+ve).
  \begin{explanation}
    If \(p\) is less than \(q\), there will be low chances
    of success so distribution will be skewed to left.\\
  \end{explanation}

  If \(p>0.5\), distribution is skewed left(-ve).
  \begin{explanation}
    If \(p\) is more than \(q\), there will be high chances
    of success so distribution will be skewed to right.\\
  \end{explanation}

  \textbf{When \(n\) is sufficiently large, binomial distribution is symmetric.}
\end{remark}

\subsection{Expectation of Binomial Distribution}
\label{ssec:expectation-of-binomial-distribution}

Given a binomial random variable \(X \sim \text{Binomial}(n,
p)\) with probability mass function:
\[
  P_X(r) = P(X = r) = \binom{n}{r} q^{n-r} p^{r}, \quad r =
  0, 1, 2, \ldots, n,
\]
where \( q = 1 - p \).

\bigskip

\begin{definition}
  The expectation \(E[X]\) is defined as:
  \[
    E[X] = \sum_{r=0}^n r \, P_X(r) = \sum_{r=0}^n r
    \binom{n}{r} q^{n-r} p^r.
  \]
\end{definition}

\bigskip

\begin{note}
  The term for \(r=0\) is zero since it is multiplied by
  \(r\), so we can start the sum from \(r=1\):
  \[
    E[X] = \sum_{r=1}^n r \binom{n}{r} q^{n-r} p^r.
  \]
\end{note}

\bigskip

Rewrite \(r \binom{n}{r}\) using the identity:
\[
  r \binom{n}{r} = n \binom{n-1}{r-1}.
\]

Therefore,
\[
  E[X] = \sum_{r=1}^n n \binom{n-1}{r-1} q^{n-r} p^r = n \sum_{r=1}^n \binom{n-1}{r-1} q^{n-r} p^r.
\]

\bigskip

Change the index of summation by letting \(k = r - 1\). When
\(r=1\), \(k=0\), and when \(r=n\), \(k=n-1\). So,
\[
  E[X] = n \sum_{k=0}^{n-1} \binom{n-1}{k} q^{n - (k+1)}
  p^{k+1} = n p \sum_{k=0}^{n-1} \binom{n-1}{k} q^{(n-1)-k}
  p^k.
\]

\bigskip

Notice that the sum is the binomial expansion of \((p +
q)^{n-1}\):
\[
  \sum_{k=0}^{n-1} \binom{n-1}{k} p^k q^{(n-1)-k} = (p + q)^{n-1} = 1^{n-1} = 1.
\]

\bigskip

Hence,
\[
  E[X] = n p \cdot 1 = n p.
\]

\bigskip

\begin{remark}
  The expectation of a binomial random variable \(X \sim \text{Binomial}(n,p)\) is
  \begin{equation}
    \label{eqn:expectation-of-binomial-distribution}
    \begin{displaystyle}
      \boxed{E[X] = n p.}
    \end{displaystyle}
  \end{equation}
\end{remark}

\subsection{Variance of Binomial Distribution}
\label{ssec:variance-of-binomial-distribution}

Recall that the
\hyperref[sec:variance-of-a-random-variable]{variance of a
random variable} \(X\) is defined as:
\[
  \operatorname{Var}(X) = E[X^2] - (E[X])^2.
\]

We have already derived that,
\footnote{See \autoref{eqn:expectation-of-binomial-distribution}}
\[
  E[X] = n p.
\]

\bigskip

Next, we compute \(E[X^2]\):
\[
  E[X^2] = \sum_{r=0}^n r^2 P_X(r) = \sum_{r=0}^n r^2 \binom{n}{r} q^{n-r} p^r,
\]
where \(q = 1-p\).

\bigskip

Rewrite \(r^2\) as \(r(r-1) + r\), so that
\begin{align*}
  E[X^2] 
&= \sum_{r=0}^n \bigl[r(r-1) + r\bigr] \binom{n}{r} q^{n-r} p^r \\
&= \sum_{r=0}^n r(r-1) \binom{n}{r} q^{n-r} p^r + \sum_{r=0}^n r \binom{n}{r} q^{n-r} p^r.
\end{align*}

\bigskip

We already know the second sum is \(E[X] = n p\). Focus on the first sum:
\[
  S = \sum_{r=0}^n r(r-1) \binom{n}{r} q^{n-r} p^r.
\]

\bigskip

Using the combinatorial identity:
\[
  r(r-1) \binom{n}{r} = n (n-1) \binom{n-2}{r-2},
\]
we get
\[
  S = n (n-1) \sum_{r=2}^n \binom{n-2}{r-2} q^{n-r} p^r.
\]

\bigskip

Change the index by letting \(k = r - 2\). When \(r=2\), \(k=0\), and when \(r=n\), \(k = n-2\):
\begin{align*}
  S 
&= n (n-1) \sum_{k=0}^{n-2} \binom{n-2}{k} q^{n - (k+2)} p^{k+2} \\
&= n (n-1) p^2 \sum_{k=0}^{n-2} \binom{n-2}{k} q^{(n-2)-k} p^k.
\end{align*}

\bigskip

Recognize the sum as the binomial expansion of \((p + q)^{n-2} = 1^{n-2} = 1\):
\[
  S = n (n-1) p^2.
\]

\bigskip

Putting it all together:
\begin{align*}
  E[X^2] &= S + E[X] \\
         &= n (n-1) p^2 + n p.
\end{align*}

\bigskip

Finally, compute the variance:
\begin{align*}
  \operatorname{Var}(X) 
&= E[X^2] - (E[X])^2 \\
&= n (n-1) p^2 + n p - (n p)^2 \\
&= n (n-1) p^2 + n p - n^2 p^2.
\end{align*}

\bigskip

Simplify:
\begin{align*}
  \operatorname{Var}(X) 
&= n p + n (n-1) p^2 - n^2 p^2 \\
&= n p + n^2 p^2 - n p^2 - n^2 p^2 \\
&= n p - n p^2 \\
&= n p (1 - p) \\
&= n p q.
\end{align*}

\bigskip

\begin{remark}
  The variance of a binomial random variable \(X \sim \text{Binomial}(n,p)\) is
  \[
    \boxed{
      \operatorname{Var}(X) = n p q,
    }
  \]
  where \(q = 1 - p\).
\end{remark}

